% ============================================================
%  Chương II – Nền tảng mật mã khóa công khai
%  ~600 từ · 2 trang
% ============================================================
\chapter{Nền tảng mật mã khóa công khai}
\label{chap:crypto}

% Ý chính: xây dựng logic dẫn truyện — symmetric (có vấn đề key exchange)
% → asymmetric (giải quyết key exchange nhưng tạo authentication problem)
% → cần PKI (giải quyết nốt).

% ----------------------------------------------------------
\section{Mật mã đối xứng và hạn chế}
\label{sec:2.1}

\subsubsection*{Nguyên lý mật mã đối xứng}
% Ý chính: cùng khóa K cho mã hóa và giải mã; ưu điểm tốc độ cao (AES, 3DES)

\subsubsection*{Bài toán phân phối khóa}
% Ý chính: n node cần O(n²) cặp khóa; cần kênh bí mật để giao khóa — vòng tròn luẩn quẩn

\subsubsection*{Giải pháp KDC và giới hạn}
% Ý chính: KDC (ví dụ Kerberos \cite{Steiner1988}) giải quyết scale
% nhưng tạo Single Point of Failure, không phù hợp môi trường internet mở

% ----------------------------------------------------------
\section{Mật mã khóa công khai}
\label{sec:2.2}

\subsubsection*{Cặp khóa bất đối xứng và trapdoor function}
% Ý chính: public key mã hóa/verify — private key giải mã/ký
% Trapdoor one-way function: dễ tính một chiều, cực khó đảo ngược

\subsubsection*{Thuật toán RSA}
% Ý chính: hard problem = phân tích thừa số nguyên tố lớn \cite{Rivest1978}
% Key gen, encrypt ($C = M^e \bmod n$), decrypt ($M = C^d \bmod n$)

\subsubsection*{Trao đổi khóa Diffie-Hellman}
% Ý chính: hai bên thoả thuận session key qua kênh công khai
% Hard problem = discrete logarithm \cite{Diffie1976}

\subsubsection*{Sơ đồ chữ ký số}
% Ý chính: $\sigma = \mathrm{Enc}_{K_{\mathrm{priv}}}(H(m))$; verify bằng public key
% Ba tính chất: authenticity, integrity, non-repudiation
% Thuật toán thực tế: RSA-PSS, ECDSA, EdDSA

% ----------------------------------------------------------
\section{Tại sao mật mã bất đối xứng chưa đủ}
\label{sec:2.3}

\subsubsection*{MITM attack trên raw public key}
% Ý chính: Mallory chặn key exchange, thay $K_{\mathrm{pub\_Bob}}$ bằng $K_{\mathrm{pub\_Mallory}}$
% → cả Alice lẫn Bob đều không phát hiện; ví dụ Bảng~\ref{tab:mitm}

\subsubsection*{Vấn đề ràng buộc danh tính}
% Ý chính: asymmetric crypto không có binding "key X thuộc về entity Y"
% → đây chính xác là lý do PKI tồn tại; bridge sang Chương~\ref{chap:pki-arch}
